\begin{theorem}
For every fixed integer $\ell\geq3$,
\[
 r(s;\ell)=\Omega\!\left(2^{(\ell-1)s/2}\right).
\]
That is, there are constants $c_\ell>0$ and $S_\ell$ such that the displayed
lower bound with factor $c_\ell$ holds for every $s\geq S_\ell$.
\end{theorem}
\begin{proof}
Fix an integer $\ell\geq3$.  For every sufficiently large integer $s$ we
have $4\leq s\leq s$, so \noderef{F2ForwardIndependentBound} gives a
digraph $D=D_2(s)$ which is loopless and $T_s$-free, with
\[
 N=2^{2s-3}-2^{s-2}
\]
vertices and
\[
 \operatorname{fwi}_s(D)\leq
 F_s:=\sum_{t=1}^{s-1}\binom{s}{t}
 2^{(s-1)(t+s)-\binom{t+1}{2}}.
\]
We now estimate this exact finite sum.  Put $L=\log_2 s$ and, in the
summand indexed by $t$, put $u=s-t$.  Then $1\leq u\leq s-1$ and
\[
 (s-1)(t+s)-\binom{t+1}{2}
 =\frac32s^2-\frac52s+1-\binom{u-1}{2}.
\]
Consequently,
\[
 F_s=2^{\frac32s^2-\frac52s+1}
 \sum_{u=1}^{s-1}\binom{s}{u}2^{-\binom{u-1}{2}}.
\]
For $u<4L$, the elementary bound $\binom{s}{u}\leq s^u$ gives
\[
 \binom{s}{u}2^{-\binom{u-1}{2}}
 \leq 2^{uL}\leq2^{4L^2}.
\]
For $u\geq4L$ and all sufficiently large $s$, we have
\[
 \binom{u-1}{2}=\frac{(u-1)(u-2)}2\geq uL,
\]
because $(u-1)(u-2)/(2u)=u/2-3/2+1/u\geq L$ in this range.
Thus the same summand is at most $s^u2^{-uL}=1$.  There are fewer
than $s$ summands, so
\[
 \sum_{u=1}^{s-1}\binom{s}{u}2^{-\binom{u-1}{2}}
 \leq s\,2^{4L^2}=2^{L+4L^2}=2^{o(s)},
\]
using $(\log_2s)^2=o(s)$.  Absorbing the constant $1$ into the
$o(s)$ term yields the required local specialization
\[
 \operatorname{fwi}_s(D)\leq
 2^{\frac32s^2-\frac52s+o(s)}.
\]

The exact value of $N$ also gives the lower bound needed below: since
$s\geq2$,
\[
 N=2^{2s-3}(1-2^{1-s})\geq2^{2s-4}>0.
\]
Choose
\[
 n=2^{\left\lfloor(s/2-4)(\ell-1)\right\rfloor}.
\]
For fixed $\ell$, this is at least $s$ for all sufficiently large $s$.
Set
\[
 Q_s=\binom ns\left(\frac{\operatorname{fwi}_s(D)}{N^s}\right)^{\ell-1}.
\]
Using $\binom ns\leq n^s$ and the two displayed estimates, an upper bound
for $Q_s$ has base-$2$ exponent at most
\begin{align*}
 &s\left\lfloor(s/2-4)(\ell-1)\right\rfloor
 +\left(-\frac12s^2+\frac32s+o(s)\right)(\ell-1)\\
 &\leq \left(-\frac52s+o(s)\right)(\ell-1)<0
\end{align*}
for all sufficiently large $s$.  Hence
\[
 Q_s<1.
\]

We may therefore apply \noderef{RandomHomomorphismColoring}.  Its
positivity hypotheses hold because $s>0$ and $\ell>0$, while the finite
order of $D$ is the positive integer $N$ above.  Its remaining graph
hypotheses hold exactly because \noderef{F2ForwardIndependentBound}
provided that $D$ is loopless and $T_s$-free, and its strict expected-count
hypothesis is the last displayed inequality $Q_s<1$.  We obtain an
$\ell$-coloring
of the complete graph on $n$ vertices with no monochromatic clique of size
$s$.

By the definition of \noderef{MulticolorRamseyNumber}, if its least
universal order were some $m\leq n$, restricting this coloring to any $m$
vertices would give an $\ell$-coloring at that order with no monochromatic
$K_s$, contradicting the universal property.  Therefore
\[
 r(s;\ell)>n.
\]
Finally,
\[
 n\geq 2^{(s/2-4)(\ell-1)-1}
 =2^{-4(\ell-1)-1}2^{(\ell-1)s/2}.
\]
Thus the fixed positive constant
$c_\ell=2^{-4(\ell-1)-1}$ and a suitable threshold $S_\ell$ give
$r(s;\ell)\geq c_\ell2^{(\ell-1)s/2}$ for every $s\geq S_\ell$.
\end{proof}
